% base/settings
% Configuraciones globales del documento.
% Modificar aquí y solo aquí los parámetros tipográficos.

%
% Márgenes de página
%
% Modificar aquí para ajustar el área de trabajo
\geometry{
  lmargin = 2.54cm,
  rmargin = 2.54cm,
  tmargin = 2.54cm,
  bmargin = 2.54cm
}

%
% Encabezado y pie de página
%
% 
\setlength{\headheight}{14.61858pt}

%
% Espaciado entre líneas
%
\doublespacing

%
% Tolerancia tipográfica (control de corte de palabras)
% Valores altos = LaTeX parte palabras menos agresivamente
\pretolerance=100
\tolerance=800


%
% Modificaciones a los estilos por defecto de book
%
% Elimina el \thispagestyle{plain} automático del \chapter
% La primera página hereda el \pagestyle activo
\patchcmd{\chapter}{\thispagestyle{plain}}{}{}{}


%
% Modificaciones en los títulos
%
% Editar el formato de títulos de capítulo
\makeatletter
\renewcommand{\@makechapterhead}[1]{% % Para capítulo numerado
  {\parindent\z@
    \hangindent=0.63cm\hangafter=1%
    \normalfont\fontsize{14}{17}\selectfont\bfseries
    \chaptername\ \thechapter.\ #1\par\nobreak}%
  \vspace{12pt}%
}
\renewcommand{\@makeschapterhead}[1]{% % Para capítulo no numerado
  {\parindent\z@\centering
    \normalfont\fontsize{14}{17}\selectfont\bfseries
    #1\par\nobreak}%
  \vspace{12pt}%
}
\makeatother


%
% Sangría y espacio entre párrafos
%
% Eliminar la sangría incial 
\setlength{\parindent}{0pt}
% Espaciado adicional entre parrafos
\setlength{\parskip}{1em}

%
% Formato de secciones en el cuerpo e índice
%
% Control del formato
% \titleformat{[nivel]}{[formato]}{[numeración]}{[espaciado entre número y título]}{[código extra]}
\titleformat{\section}{\normalsize\bfseries}{\thesection}{1em}{}
\titleformat{\subsection}{\normalsize\bfseries}{\thesubsection}{1em}{}
\titleformat{\subsubsection}{\normalsize\bfseries}{\thesubsubsection}{1em}{}
% Profundidad de numeración y tabla de contenidos
\setcounter{secnumdepth}{3}
\setcounter{tocdepth}{3}
% Control del espaciado
% \titlespacing{[nivel]}{[sangría izquierda]}{[espacio antes del título]}{[espacio después del título]}
% Formato de la entrada de capítulo en el índice
% {nivel}[sangría]{código arriba}{etiqueta numerada}{etiqueta sin número}{relleno + página}
\titlecontents{chapter}[0pt]
  {\bfseries}
  {\chaptername\ \thecontentslabel.\enspace}
  {}
  {\hfill\contentspage}
\titlecontents{section}[1.5em]
  {\bfseries}
  {\thecontentslabel\enspace}
  {}
  {\hfill\contentspage}
\titlecontents{subsection}[3.8em]
  {}
  {\thecontentslabel\enspace}
  {}
  {\hfill\contentspage}


%
% Formato de tablas y figuras en el cuerpo
%
% Aplica doble espaciado a ambas porque por defecto es 1
\captionsetup{font={stretch=2}}
% Redefine cómo iría el caption #1: Tabla 1.1, #2: :, #3: Título
\DeclareCaptionFormat{tablacaption}{\textbf{#1}\\\textit{#3}}
% Redefine el formato de tabla: labelsep suele ser ":" y singlelinecheck suele
% poner centrado si la linea cabe
\captionsetup[table]{
  format=tablacaption,
  labelsep=none,
  justification=justified,
  singlelinecheck=false
}
% Redefine el formato de figura: solo pone Figura 1.1 como negrita
\captionsetup[figure]{labelfont=bf}



%
% Configuración de display math
%
% Permite que ecuaciones largas corten entre páginas
% El número es el nivel de permisividad
\allowdisplaybreaks[1]



